% Root and B literal audits passed; independent equal-loop and isolated-image/lift gates passed.
\subsection{Disconnected six-state graphs with an eight-element source}

\begin{lemma}[Low homology under an action-image quotient]
\label{C:graph-action-image-low-homology}
Let a graph $\Gamma$ carry a cellular right action of a monoid $M$
which factors through a surjection $M\to H$.  The actual map of
action-bar cofibers
\[
 X_M=\operatorname{hocofib}(\Gamma_{hM}\to BM)
       \longrightarrow
 X_H=\operatorname{hocofib}(\Gamma_{hH}\to BH)
\]
is an isomorphism on $H_1$ and a surjection on $H_2$, integrally
and over every field.  The graph need not be connected.
\end{lemma}
\begin{proof}
Use the augmented graph complex
$A_\Gamma=[C_1(\Gamma)\to\widetilde C_0(\Gamma)]$,
in degrees one and zero.  A natural reduced chain model for the
cofiber is the suspension of the total two-sided monoid bar on
$A_\Gamma$.  Before that suspension, the quotient map is literally
the identity in total degree zero and is onto in total degree one.
The images of the two differentials into degree zero coincide.
Consequently it is an isomorphism on degree-zero homology.
Every degree-one target cycle lifts to a source chain; its boundary
is zero since degree zero is unchanged.  This proves surjectivity
on degree-one homology, hence the asserted shifted statements.
\end{proof}

\begin{lemma}[A graph quotient with no cycle coinvariants in its kernel]
\label{C:graph-cycle-kernel-quotient}
Suppose an equivariant graph quotient $\Gamma\to\Gamma'$ is
bijective on components and onto on first homology.  Write $L$
for the kernel of its first-homology map over a field $k$.
If $L\otimes_{kM}k=0$, its action-bar cofiber map induces an
isomorphism on $H_2(-;k)$.
\end{lemma}
\begin{proof}
The actual cellular chain maps used here are degreewise
surjective.  Their first-homology surjectivity and component
isomorphism imply that the chain kernel has homology only $L$
in degree one.  Graphs have no cellular two-boundaries, so the
literal inclusion of kernel cycles is a natural equivariant
quasi-isomorphism from $L$ in that degree.  For an arbitrary
cellular comparison with the stated homology properties, use its
derived fiber and the same mapping-cone homology sequence instead.
After taking the free monoid bar and suspending, the comparison
fiber has homology $L\otimes_{kM}k$ in degree two and zero in
degree one.  The homology exact sequence gives the assertion.
\end{proof}

\begin{proposition}[Two bounded disconnected image controls]
\label{C:right-zero-six-disconnected-controls}
Consider the following two six-state right-zero action graphs,
up to relabeling states and interchanging $p,q$ (which reverses
all edge orientations).
\begin{enumerate}
\item There is a loop at $0$, a directed edge $0\to1$, and a
separate loop at $2$.  Every right actor preserves the functional
which is one on both loops.
\item There is a loop at each of $0,1$, an edge $0\to1$, and an
isolated vertex $5$.  The restriction of the right-action image
to the five-state connected component is the complete five-element
mixed image from Theorem~\ref{C:right-zero-five-connected}.
\end{enumerate}
If $|M|\leq8$, neither collage has type $\mathbb CP^2$.
\end{proposition}
\begin{proof}
For the first graph, the two left retractions on the state list
$(0,1,2,\ell_0,\ell_2,a)$ are
\[
 P=(0,1,2,0,2,0),\qquad Q=(0,1,2,0,2,1).
\]
There are exactly $55$ commuting state maps.  Exactly ten preserve
the displayed loop functional, and these form a monoid.  Among its
$1024$ subsets, exactly $56$ contain the identity, are closed, and
have cardinality at most eight.  Their literal collages split as
follows: seven have nonzero $H_1(\mathbb F_2)$; three have
$\dim H_2(\mathbb F_2)=2$; and forty-six have one-dimensional
$H^2(\mathbb F_2)$ with an explicit zero-square primitive.
No image has zero second homology.

The producer record is
\path{code/outputs/a_rightzero_six_loop_tail_images.json}; only
its first, equal-loop hull is used here.  The standalone verifier
\path{code/d_verify_rightzero_six_equal_loop.py} independently
enumerates all $6^6$ state maps and all $1024$ subsets, rebuilds
every literal category, and verifies its binary ranks and actual
cup equations.  Its complete record
\path{code/outputs/d_rightzero_six_equal_loop_independently_verified.json}
checks $23169$ typed associativity equations, $11659$ cocycle
equations, and $58771$ primitive equations.
Lemma~\ref{C:graph-action-image-low-homology} transfers these
obstructions to every source with that image.  In the square case,
the target's nonzero one-dimensional second homology and the
surjection from the hypothetical projective plane make pullback
an isomorphism on second cohomology; this nonzero-target condition
is essential.

For the second graph put
\[
 P=(0,1,0,1,0,5),\qquad Q=(0,1,0,1,1,5).
\]
A right actor restricts to one of the five prescribed maps on
$0,\ldots,4$ and sends $5$ to one of the three core vertices
$0,1,5$.  These are precisely fifteen maps, closed under
composition.  Exhausting their subsets of cardinality at most
eight whose restriction contains all five component maps gives
the following $93$ literal image monoids:
\[
\begin{array}{c|rrrr}
\text{image order}&5&6&7&8\\\hline
H_1(\mathbb F_2)\ne0&1&0&0&0\\
\dim H_2(\mathbb F_2)=2&1&0&0&0\\
\dim H_2(\mathbb F_2)=1,\ \text{square zero}&0&10&12&6\\
H_2(\mathbb F_2)=0&0&0&19&44
\end{array}
\]
The first three rows are excluded by the same actual comparison.
An order-eight image must be the source itself, so its zero
second homology also excludes it.  An order-seven image leaves
only the source itself or an eight-element source with one
two-element quotient fiber.

If $H^2(BM;\mathbb Q)=0$, the collage's degree-two classes are
relative to its two vertex spaces, and the suspension quotient
kills their squares.  Otherwise every order-eight source occurs
in the complete $516$-table list proved in
Theorem~\ref{C:eight-act-cones}.  No published-census completeness
is needed for that list.  For each such table, try every pair
of elements as the sole nontrivial congruence class.  There are
$989$ actual one-pair congruences.  Construct the seven-element
quotient and try every literal isomorphism to each of the nineteen
remaining image monoids.  This yields $660$ distinct right acts
on the fixed six states.  Every one of their actual seventeen-arrow
collages has $H_2(\mathbb F_2)=0$.

The complete producers and records are
\path{code/a_rightzero_six_isolated_images.py},
\path{code/a_rightzero_six_isolated_inflations.py}, and their
same-prefix files in \path{code/outputs/}.
The latter record stores the source index, unique quotient pair,
literal quotient map, full category table, and homology result
for every act.  It links the original $516$-table source and the
$93$-image record by their exact hashes.
The standalone verifier
\path{code/b_verify_rightzero_six_isolated.py} imports no producer
or matrix-library routines.  It reconstructs the fifteen maps
from all $6^6$ state maps and exhausts all $16384$ subsets containing
the identity.  It independently checks all $516$ source tables,
all $989$ pair congruences, and all $19\cdot720$ target relabelings;
their literal quotient lifts reproduce exactly the $660$ acts.
All $753$ image and lifted categories have their typed laws and
normalized bar equations checked, including $89048$ supplied
cup-primitive equations.  Its complete record is
\path{code/outputs/b_rightzero_six_isolated_independently_verified.json}.
Thus the nineteen nonfaithful-source cases are excluded by a
complete small lift calculation, not by transferring zero second
homology from an action image.
\end{proof}

\begin{theorem}[Six right-zero states with a source of order at most eight]
\label{C:right-zero-six-small-sources}
Let $R=\{1,p,q\}$ be the identity-adjoined right-zero two-band.
For $|M|\leq8$ and $|S|\leq6$, an $(R,M)$-biset collage cannot
have type $\mathbb CP^2$.  In particular the seventeen-arrow
profile $(|M|,|R|,|S|)=(8,3,6)$ is excluded.
\end{theorem}
\begin{proof}
The connected case is Theorem~\ref{C:right-zero-six-connected}.
Let $\Gamma$ be disconnected and put $T=\pi_0\Gamma$.
The ordinary $T$-act cone has zero rational degree-two products
by Lemma~\ref{C:nine-unbounded-act-cups}.  If a hypothetical
projective-plane generator pulled back from that cone, its square
would therefore be zero.  By the actual augmented-graph exact
sequence of Lemma~\ref{C:graph-component-eigenvalues}, its image
is consequently a nonzero invariant functional
\[
 \lambda:H_1(\Gamma;\mathbb Q)\longrightarrow\mathbb Q.
\]
This assertion does not identify the entire second cohomology
with cycle coinvariants when the graph is disconnected.

A forest is an actual ordinary-act reduction.  If all edges are
loops, the vertex retraction is an ordinary-act homotopy retract
with suspension relative quotient.  Its second cohomology either
spans that of the collage, or every degree-two class is relative;
both possibilities give zero square.  We may assume there is a
nonloop edge and a cycle.

With four vertices and at most two edges, a unique loop forces
its vertex to be fixed by every actor unless its cycle
coinvariants vanish.  The only other cyclic possibility is a
two-edge circle and isolated vertices.  Nonzero $\lambda$ forces
degree $+1$ on the circle and preservation of its component.
The degree-three circle map fixes its two vertices and commutes
with identity or half-rotation; extend it by the identity on all
isolated vertices.  It remains equivariant even when an isolated
vertex enters the circle.  The component eigenvalue lemma kills
the square.  More vertices leave no further nonforest case.

Thus there are three vertices, three edges, and two components.
If all cycles lie in a single component, either its vertex is
fixed, or it has two vertices and three edges, with the remaining
vertex isolated.  Call that connected component $C$.  Every actor
must preserve $C$: sending it to the isolated vertex annihilates
the nonzero functional $\lambda$.

Apply the literal connected five-state analysis to the action on
$C$.  Its reductions use a nonzero invariant cycle functional and
absence of a common fixed core vertex.  With three connecting
edges, a swap kills the parallel-edge difference.  With two
connecting edges and one loop, a swap kills the loop difference.
The corresponding equivariant quotients preserve the isolated
vertex and have zero coinvariants in their cycle kernel, so
Lemma~\ref{C:graph-cycle-kernel-quotient} gives a second-homology
isomorphism.  Their circle power maps fix every graph vertex and
therefore extend equivariantly over the isolated vertex.

For one connecting edge and two loops, both loops must be based
at different endpoints; otherwise their common loop vertex is
fixed.  The nonzero functional and absence of a fixed core vertex
force both constant-vertex types and equal nonzero loop values.
If all constants have the same collapse choice, the strict
degree-two maps in Theorem~\ref{C:right-zero-five-connected}
fix every vertex and extend over the isolated vertex.  The
component eigenvalue lemma again applies.  The only remaining
case is the complete mixed five-element component image, which
is precisely the second control of the preceding proposition.
No coinvariant-rank assumption about the disconnected collage
has been inserted in this reduction.

Suppose instead that both components contain a cycle.  One
component has two vertices and the other is a single vertex.
If the two-vertex component is a two-edge circle, the other is
a loop.  A nonzero value on that unique loop forces its vertex
to be fixed by every actor.  Otherwise $\lambda$ is nonzero on
the circle and zero on the loop.  No actor can send the circle
to the loop component, and its degree on the circle must be
$+1$.  Apply the degree-three map to both the circle and the
loop.  It fixes all vertices and commutes also with an actor
which collapses the loop to a circle vertex.  Thus it is a strict
scalar graph map and the component eigenvalue lemma applies.

The remaining two-vertex component is a loop with a connecting
tail, while the singleton component also carries one loop.
If the two loop values of $\lambda$ are unequal, a loop with
nonzero value must be sent to itself by every actor, forcing
its core vertex to be fixed.  If they are equal, they are both
nonzero, and every actor belongs to the equal-loop hull in the
first control above.  This excludes the last case.
Finally, a tree component and a singleton component carrying
both loops also force the latter vertex to be fixed: sending
that component into the tree annihilates $\lambda$.

This exhausts the disconnected graphs and proves the theorem.
\end{proof}
